\section{Lema tècnic}
En aquesta secció ens preocuparem de la convergència de la sèrie
\[
\sum_{p\text{ primer}}\frac{k-\omega_f(p)}{p}
\]
on $f:=f_1\cdots f_k$ per $f_1,\ldots,f_k\in\ZZ[X]$ irreductibles diferents amb el coeficient de l'exponent més elevat positiu.

Vegem primer que ens podem reduir a estudiar la convergència de la mateixa sèrie però prenent $f$ com un irreductible.
\begin{lem}
    Per a tots els primers $p$ llevat d'un nombre finit tenim
    \[
    \omega_f(p)=\omega_{f_1}(p)+\cdots +\omega_{f_k}(p).
    \]
\end{lem}
\begin{proof}
    És clar que $\omega_f(p)\leq \omega_{f_1}(p)+\cdots +\omega_{f_k}(p)$. Només cal veure que $f_i$ i $f_j$ per $i\neq j$ no tenen zeros en comú per quasi totes les $p$ ---és a dir, per totes llevat d'un nombre finit---. Pel lema de Gauss $f_i$ i $f_j$ són irreductibles de $\QQ[X]$ i, per tant, $\gcd(f_i,f_j)=1$ ja que són diferents. Sabem que $\QQ[X]$ és un domini Euclidià, en particular tenim la identitat de Bézout,
    \[
    uf_i+vf_j=1
    \]
    per $u,v\in \QQ[X]$. Sigui $d\in \ZZ_{\geq 1}$ el mínim enter tal que $du, dv\in\mathbb Z[X]$. Prenem $g:=du$ i $h:=dv$. Aleshores,
    \[
    gf_i+hf_j=d.
    \]
    Suposem que existeix un zero comú de $f_i$ i $f_j$ en $\FF_p$ per $p$ un primer. Llavors $d\equiv 0\pmod{p}$. Però aquesta igualtat només pot passar per un nombre finit de primers ja que $\mathbb Z$ és un domini de factorització única. En particular, essent més grollers ja que només ens importa la convergència de certa sèrie, per tot primer $q>d$ tenim la igualtat de l'enunciat.
\end{proof}

Aquest lema ens permet expressar la sèrie de la següent forma\footnote{Per dos nombres enters $x,y\in\ZZ_{\geq 1}$ escrivim $\sum_{p\leq x}$, $\sum_{y<p\leq x}$ per denotar la suma sobre tots els primers $p\leq x$ o tots els primers $y<p\leq x$ respectivament.}
\[
\sum_{p\leq x}\frac{\omega_f(p)-k}{p}=\sum_{n=1}^k\sum_{d<p\leq x}\frac{\omega_{f_i}-1}{p}+\sum_{n=1}^k\sum_{p\leq d}\frac{\omega_{f_i}(p)-1}{p}
\]
on $d\in\ZZ_{\geq 1}$ és tal que per tot primer $p>d$ tenim la igualtat del lema anterior i $x>d$ és qualsevol nombre enter.

A més a més, sigui $a_i$ el coeficient d'exponent més alt del polinomi $f_i$ i siguin $(p_j^{(i)})_{j=1}^{n_i}$ els primers de la factorització de $a_i$ per $i=1,\ldots, k$. Si considerem $d':=\max(d, \max_{i,j}p_j^{(i)})$ aleshores podem reduir-nos encara més ja que el polinomi $f_i$ dins $\FF_q$ per $q> d'$ té el mateix nombre d'arrels que el polinomi mònic $g_i:=a_i^{-1}f_i\in\FF_q[X]$. En resum, per $x>d'$, tenim
\[
\sum_{p\leq x}\frac{\omega_f(p)-k}{p}=\sum_{n=1}^k\sum_{d'<p\leq x}\frac{\omega_{f_i}-1}{p}+\sum_{n=1}^k\sum_{p\leq d'}\frac{\omega_{f_i}(p)-1}{p}.
\]
Només ens cal estudiar la convergència de la sèrie 
\[
\sum_{p\text{ primer}}\frac{\omega_g(p)-1}{p}
\]
per $g\in\ZZ[X]$ irreductible i mònic.

\subsection{Teoremes algebraics sense demostració}

Per tal de reduir al màxim la demostració, enunciarem uns quants resultats de teoria de nombres algebraica citant, quan sigui necessari, a on es pot trobar la demostració d'aquests. Els resultats següents ens ajudaran a estudiar el nombre de solucions d'un polinomi mònic irreductible a $\FF_p$.

D'ara en endavant tindrem $K/\QQ$ una extensió finita i $\ZZ_K$ l'anell dels enters algebraics.
\begin{thm}[de l'element primitiu]
    Existeix $\theta\in \ZZ_k$ tal que $K=\QQ(\theta)$.
\end{thm}
\begin{thm}[Teorema 5.6 i 5.3d de \cite{ANT}]
    Tot ideal de $\ZZ_K$ descompon en producte d'ideals primers únics llevat d'ordre. A més a més, tot ideal primer és maximal.
\end{thm}

Sigui $p\in\ZZ$ un primer considerem l'ideal $p\ZZ_K$. Aquest descompon en producte d'ideals primers. Siguin $P_1,\ldots,P_n\subseteq \ZZ_K$ els ideals primers i $e_1,\ldots,e_n$ els exponents tals que
\[
p\ZZ_K=P_1^{e_1}\cdots P_n^{e_n}.
\]
Com que els ideals primers de $\ZZ_K$ són maximals tenim que $\ZZ_K/P_i$ és un cos de característica $p$ ---teorema 27 de \cite{NF}---. En particular, $\#(\ZZ_K/P_i)=p^{f_i}$ per un $f_i\in\ZZ_{\geq 1}$ determinat que anomenarem \emph{grau d'inèrcia} de $p$ dins $P_i$. 

D'ara en endavant denotarem per $N(P_i):=\#(\ZZ_K/P_i)$ i en direm la \emph{norma} de $P_i$ ---de fet aquesta definició la tenim per a tot ideal $0\neq I\subseteq \ZZ_K$---.

\begin{thm}[Criteri de factorització de Dedekind]
    Sigui $K=\QQ(\theta)$ per $\theta\in\ZZ_K$ i sigui $p\in \ZZ$ un primer tal que $p\ZZ_K$ factoritza de la següent forma
    \[
    p\ZZ_K=P_1^{e_1}\cdots P_n^{e_n},
    \]
    on $P_1,\ldots, P_n$ són ideals primers. Si $p\nmid [\ZZ_K:\ZZ[\theta]]$, aleshores existeix una factorització 
    \[
    \operatorname{irr}(\theta, \QQ)\equiv g_1^{e_1}\cdots g_n^{e_n}\pmod{p}
    \]
    on els $g_i$ són irreductibles a $\FF_p[X]$ amb $\deg g_i=f_i$, el grau d'inèrcia de $p$ dins $P_i$, per a tot $i=1,\ldots, n$.
\end{thm}
Aquest teorema ens diu que $[K:\QQ]=\sum_{i=1}^ke_if_i$. A més, $\operatorname{irr}(\theta, \QQ)$ té arrels a $\FF_p$ si i només si $g_i=X-a$ per un $a\in \FF_p$ i un $i$ concret i, per tant, $f_i=1$ d'on obtenim que la norma de $P_i$ és $p$. Aquest anàlisi ens ha demostrat el següent coro\lgem ari,
\begin{coro}\label{coro}
Sigui $g\in \ZZ[X]$ irreductible i mònic amb arrel $\theta$, en alguna clausura algebraica. Considerem l'extensió $K:=\QQ(\theta)$. Per quasi tots els primers $p$; és a dir, per a tots menys un nombre finit, el nombre $\omega_g(p)$ equival al nombre d'ideals de norma $p$ de la descomposició en ideals primers de $p\ZZ_K$.  
\end{coro}

\subsection{Teorema de Mertens}

Només ens queda enunciar dos últims teoremes, les demostracions dels quals es troben a molts llocs, per poder demostrar el lema tècnic, i.e., la convergència de 
\[
\sum_{p\text{ primer}}\frac{\omega_g(p)-1}{p}
\]
per $g\in\ZZ[X]$ irreductible i mònic.

Aquests dos teoremes es podrien muntar en un de sol, per completesa els escriurem separats,
\begin{thm}[de Mertens]
Sigui $x\in\mathbb Z_{\geq 2}$.
    \[
    \sum_{p\leq x}\frac{1}{p}=\log\log x+B+O\left(\frac{1}{\log x}\right),
    \]
    on $B$ és la Meissel.Mertens constant.
\end{thm}
\begin{thm}[de Mertens per a cossos de nombres]
    Sigui $K/\QQ$ una extensió finita, aleshores exiteix una constant $C$ tal que 
    \[
    \sum_{N(P)\leq x}\frac{1}{N(P)}=\log\log x+C+O\left(\frac{1}{\log x}\right),
    \]
    on la suma recorre tots els ideals primers de $\ZZ_K$ de norma menor o igual a $x\in\ZZ_{\geq 1}$.
\end{thm}

Amb aquests dos i tots els resultats previs podem demostrar el següent
\begin{lem}[tècnic]
    Sigui $g\in\ZZ[X]$ mònic i irreductible. Aleshores
    \[
    \sum_{p\text{ primer}}\frac{\omega_g(p)-1}{p}
    \]
    convergeix.
\end{lem}
\begin{proof}
    Sigui $\theta$ una arrel de $g$ en algun cos de descomposició i considerem l'extensió finita $K=\QQ(\theta)$. Pel coro\lgem ari \ref{coro} tenim
    \[
    \sum_{p\leq x}\frac{\omega_g(p)}{p}=\sum_{N<N(P)\leq x}\frac{1}{N(P)}+\sum_{p\leq N}\frac{\omega_g(p)}{p}.
    \]
    Aleshores, pels teoremes de Mertens,
    \begin{align*}
        \sum_{p\text{ primer}}\frac{\omega_g(p)-1}{p}&=\sum_{N<N(P)\leq x}\frac{1}{N(P)}-\sum_{p\leq x}\frac{1}{p}+\sum_{p\leq N}\frac{\omega_g(p)}{p}\\
        &=\sum_{p\leq N}\frac{\omega_g(p)}{p}-B+C+O\left(\frac{1}{\log x}\right),
    \end{align*}
    per $B$ i $C$ les constants enunciades al teoremes. Si fem tendir $x$ a $+\infty$ obtenim el resultat desitjat.
\end{proof}

% \begin{thm}[Teorema 5.14c de \cite{ANT}]
    
% \end{thm}
